% =====================================================================
%  CHAPITRE 9 — GRAMMAIRE A2 : LES TEMPS DU VERBE
% =====================================================================
\chapter{Les temps : parler du passé, du présent et de l'avenir}

Le système verbal anglais est \textbf{plus simple que le français} (pas de passé simple
littéraire, pas de subjonctif obligatoire, pas d'imparfait distinct), mais il fait une
distinction que le français ignore : \textbf{l'action est-elle terminée ou encore ouverte ?}
C'est tout le secret du \en{present perfect}.

\section{La carte générale des temps}
\begin{tableau}[Les 8 temps que vous utiliserez réellement]
\footnotesize
\begin{longtable}{@{}p{2.9cm}p{3.6cm}p{4.2cm}p{3.6cm}@{}}
\entete{Temps} & \entete{Construction} & \entete{Emploi principal} & \entete{Exemple} \tabularnewline
\endfirsthead
\entete{Temps} & \entete{Construction} & \entete{Emploi} & \entete{Exemple} \tabularnewline
\endhead
\hline
\textbf{Présent simple} & base + \en{s} & habitude, vérité générale & \en{I work here.} \tabularnewline
\textbf{Présent continu} & \en{be} + \en{-ing} & action en cours maintenant & \en{I am working.} \tabularnewline
\textbf{Passé simple} & \en{-ed} / irrégulier & action terminée à un moment précis du passé & \en{I worked yesterday.} \tabularnewline
\textbf{Passé continu} & \en{was/were} + \en{-ing} & décor, action en cours à un moment du passé & \en{I was working at 8 pm.} \tabularnewline
\textbf{\en{Present perfect}} & \en{have} + participe passé & lien passé-présent, expérience, résultat & \en{I have worked here for 3 years.} \tabularnewline
\textbf{\en{Perf. continuous}} & \en{have been} + \en{-ing} & durée d'une activité encore en cours & \en{I have been working since 7 am.} \tabularnewline
\textbf{\en{Past perfect}} & \en{had} + participe passé & passé antérieur à un autre passé & \en{I had already eaten when he called.} \tabularnewline
\textbf{Futur} & \en{will} / \en{going to} / présent continu & décision, intention, plan & \en{I'll call you. I'm going to study.} \tabularnewline
\end{longtable}
\end{tableau}

\section{Unité 1 --- Le présent continu (\en{be} + \en{-ing})}

\begin{regle}[Formation]
\en{am / is / are} + verbe en \en{-ing} : \en{I'm working, she's reading, they're coming}.\\
Négation : \en{I'm not working.} \quad Question : \en{Are you working?} \en{What are you doing?}
\end{regle}

\begin{tableau}[Orthographe du \en{-ing}]
\begin{tabularx}{\linewidth}{@{}p{5.4cm}X@{}}
\entete{Cas} & \entete{Exemples} \tabularnewline
\hline
Règle générale : + \en{ing} & \en{work} \ra \en{working}, \en{play} \ra \en{playing}, \en{read} \ra \en{reading} \tabularnewline
Verbe se terminant par \en{e} muet : \en{e} \ra \en{ing} & \en{make} \ra \en{making}, \en{write} \ra \en{writing}, \en{come} \ra \en{coming}, \en{live} \ra \en{living} \tabularnewline
\en{-ee}, \en{-ye}, \en{-oe} : on garde tout & \en{see} \ra \en{seeing}, \en{agree} \ra \en{agreeing}, \en{dye} \ra \en{dyeing} \tabularnewline
Une syllabe, voyelle + consonne : on double & \en{sit} \ra \en{sitting}, \en{run} \ra \en{running}, \en{get} \ra \en{getting}, \en{stop} \ra \en{stopping}, \en{begin} \ra \en{beginning} \tabularnewline
\en{-ie} \ra \en{-ying} & \en{die} \ra \en{dying}, \en{lie} \ra \en{lying}, \en{tie} \ra \en{tying} \tabularnewline
\end{tabularx}
\end{tableau}

\subsection{Quand l'employer ?}
\begin{enumerate}
  \item \textbf{Action en cours au moment où je parle} : \en{Be quiet, I'm working.}
  \item \textbf{Situation temporaire} : \en{I'm living with my brother this month.}
  \item \textbf{Tendance actuelle, changement} : \en{Prices are rising. My English is improving.}
  \item \textbf{Projet futur déjà décidé} : \en{I'm travelling to Dakar tomorrow.}
  \item \textbf{Action répétée avec \en{always}} (agacement) : \en{He's always losing his keys!}
\end{enumerate}

\begin{tableau}[Présent simple contre présent continu]
\begin{tabularx}{\linewidth}{@{}p{5.4cm}p{5.4cm}X@{}}
\entete{Simple (habituel, permanent)} & \entete{Continu (maintenant, temporaire)} & \entete{Français} \tabularnewline
\hline
\en{I work in a bank.} & \en{I'm working on a report.} & Je travaille dans une banque. / Je suis en train de travailler sur un rapport. \tabularnewline
\en{She drinks tea every morning.} & \en{She's drinking tea now.} & Elle boit du thé chaque matin. / Elle boit du thé maintenant. \tabularnewline
\en{It rains a lot in August.} & \en{It's raining.} & Il pleut beaucoup en août. / Il pleut (là). \tabularnewline
\en{They live in Thiès.} & \en{They're staying at a hotel.} & Ils habitent Thiès. / Ils logent à l'hôtel. \tabularnewline
\end{tabularx}
\end{tableau}

\begin{astuce}[Les mots qui déclenchent le continu]
\en{now, right now, at the moment, currently, today, this week, these days, look!, listen!,
still, temporarily} --- dès que vous les voyez, pensez \en{-ing}.
\end{astuce}

\begin{piege}[Le français « je travaille » est ambigu, l'anglais non]
« Je travaille dans une banque » = \en{I work in a bank} (permanent).\\
« Je travaille en ce moment » = \en{I'm working right now} (en cours).\\
Si vous dites \en{I'm working in a bank}, un anglophone comprend « temporairement, cette
période-ci ». Choisissez toujours selon le sens, jamais selon la traduction mot à mot.
\end{piege}

\begin{exo}[1.1 --- Simple ou continu ?]
\begin{enumerate}
  \item Look! It (rain) \trou[2cm].
  \item She usually (take) \trou[2cm] the bus, but today she (walk) \trou[2cm].
  \item I (not / understand) \trou[2.5cm] this exercise.
  \item They (build) \trou[2cm] a new bridge at the moment.
  \item He (know) \trou[2cm] the answer.
  \item We (study) \trou[2cm] English every day.
  \item Shh! The baby (sleep) \trou[2cm].
  \item My brother (be) \trou[2cm] a teacher, but this summer he (work) \trou[2cm] in a shop.
\end{enumerate}
\end{exo}
\begin{corrige}[1.1]
1. \en{is raining} 2. \en{takes} --- \en{is walking} 3. \en{don't understand}
4. \en{are building} 5. \en{knows} 6. \en{study} 7. \en{is sleeping}
8. \en{is} --- \en{is working}
\end{corrige}

\section{Unité 2 --- Le passé simple (\en{past simple})}

\begin{regle}[Formation]
\textbf{Verbes réguliers :} base + \en{-ed} --- \en{I worked, she played, we watched}.\\
\textbf{Verbes irréguliers :} 2\textsuperscript{e} forme à apprendre par c\oe ur
(\en{go} \ra \en{went}, \en{see} \ra \en{saw}, \en{eat} \ra \en{ate}). Voir l'annexe A.\\
\textbf{\en{be} :} \en{was} (\en{I/he/she/it}) / \en{were} (\en{you/we/they}).\\
\textbf{Négation :} \en{didn't} + \textbf{base} --- \en{I didn't go} (jamais \en{*didn't went}).\\
\textbf{Question :} \en{did} + sujet + \textbf{base} --- \en{Did you see him?}
\end{regle}

\begin{tableau}[Prononciation du \en{-ed} : trois sons, pas un]
\begin{tabularx}{\linewidth}{@{}p{2.4cm}p{4.4cm}X@{}}
\entete{Son} & \entete{Après} & \entete{Exemples} \tabularnewline
\hline
/t/ & consonne sourde (\en{p, k, f, s, sh, ch, x}) & \en{worked} \pron{oueu-kt}, \en{stopped} \pron{stopt}, \en{watched} \pron{ouo-tcht} \tabularnewline
/d/ & consonne sonore ou voyelle & \en{played} \pron{plé-id}, \en{called} \pron{kold}, \en{lived} \pron{livd} \tabularnewline
/\ipaiot d/ & \en{t} ou \en{d} final & \en{wanted} \pron{ouon-tid}, \en{needed} \pron{nii-did}, \en{decided} \pron{di-saï-did} \tabularnewline
\end{tabularx}
\end{tableau}

\subsection{Quand employer le passé simple ?}
\textbf{Une action terminée, à un moment déterminé du passé.} Les mots-signaux sont vos
amis : \en{yesterday, last night / week / year, two days ago, in 2010, when I was a child,
then, at that time}.

\ex{I visited Dakar last month. She didn't come yesterday. Did you call him?}{J'ai visité Dakar le mois dernier. Elle n'est pas venue hier. L'as-tu appelé ?}

\begin{piege}[\en{*I have seen him yesterday}]
Dès qu'un \textbf{moment passé précis} est mentionné (\en{yesterday}, \en{last week},
\en{in 2015}, \en{two hours ago}), c'est le \textbf{passé simple} obligatoire. Le
\en{present perfect} est alors une faute grave --- et les francophones la font pendant des
années. Comparez :\\
\en{I saw him yesterday.} \ok \qquad \en{*I have seen him yesterday.} \ko\\
\en{I have seen him.} \ok (sans date : le lien avec le présent compte)
\end{piege}

\begin{piege}[\en{was} / \en{were} et le subjonctif \en{were}]
\en{I was tired, they were happy}. Mais dans une hypothèse irréelle, on emploie
\en{were} pour toutes les personnes : \en{If I \textbf{were} you\ldots} (Si j'étais toi\ldots),
\en{I wish I \textbf{were} taller}. Forme familière acceptée : \en{If I was you\ldots}
\end{piege}

\begin{exo}[2.1 --- Mettez au passé simple]
\begin{enumerate}
  \item \en{go, see, eat, have, do, make, take, come, say, tell}
  \item \en{buy, bring, think, teach, catch, meet, feel, leave, find, know}
  \item \en{work, play, watch, study, stop, live, want, need, call, try}
\end{enumerate}
\end{exo}
\begin{corrige}[2.1]
1. \en{went, saw, ate, had, did, made, took, came, said, told}
2. \en{bought, brought, thought, taught, caught, met, felt, left, found, knew}
3. \en{worked, played, watched, studied, stopped, lived, wanted, needed, called, tried}
\end{corrige}

\begin{exo}[2.2 --- Traduisez]
\begin{enumerate}
  \item Je ne suis pas venu hier. \item Est-ce que tu as mangé ?
  \item Elle a acheté une voiture l'année dernière. \item Nous étions très fatigués.
  \item Ils ne savaient pas la réponse. \item Quand je vivais à Dakar, je travaillais dans un magasin.
\end{enumerate}
\end{exo}
\begin{corrige}[2.2]
1. \en{I didn't come yesterday.} 2. \en{Did you eat?} 3. \en{She bought a car last year.}
4. \en{We were very tired.} 5. \en{They didn't know the answer.}
6. \en{When I lived in Dakar, I worked in a shop.}
\end{corrige}

\section{Unité 3 --- Le passé continu et le récit au passé}

\begin{regle}
\en{was / were} + \en{-ing}. Emplois :
\textbf{(1)} décor d'une histoire : \en{The sun was shining and the birds were singing.}
\textbf{(2)} action en cours interrompue par une autre :
\en{I was sleeping when the phone rang.}
\textbf{(3)} deux actions simultanées : \en{While I was cooking, he was watching TV.}
\end{regle}

\begin{tableau}[\en{when} contre \en{while}]
\begin{tabularx}{\linewidth}{@{}p{5.6cm}X@{}}
\entete{Structure} & \entete{Exemple} \tabularnewline
\hline
\en{was/were + ing} + \en{when} + passé simple & \en{I was reading when she arrived.} (action longue interrompue) \tabularnewline
\en{passé simple} + \en{while} + \en{was/were + ing} & \en{She arrived while I was reading.} \tabularnewline
\en{while} + \en{ing} (sujet commun) & \en{While walking, I saw him.} \tabularnewline
Deux passés continus & \en{While I was working, they were playing.} \tabularnewline
\end{tabularx}
\end{tableau}

\begin{astuce}[Raconter une histoire : le cocktail des trois passés]
Un bon récit anglais mélange trois temps, toujours dans le même ordre logique :\\
\textbf{1. Passé continu} pour le décor : \en{It was raining and I was waiting for the bus.}\\
\textbf{2. Passé simple} pour les événements : \en{Suddenly a man came and spoke to me.}\\
\textbf{3. \en{past perfect}} pour l'antériorité : \en{I realised I had met him before.}\\
Maîtrisez ce trio et vos récits sonneront immédiatement « anglais ».
\end{astuce}

\begin{exo}[3.1 --- Passé simple ou passé continu ?]
\begin{enumerate}
  \item I (walk) \trou[2cm] home when I (see) \trou[2cm] an accident.
  \item While she (cook) \trou[2cm], the lights (go) \trou[2cm] out.
  \item What (you / do) \trou[2.5cm] at 9 pm yesterday?
  \item He (not / listen) \trou[2.5cm], so he (not / understand) \trou[2.5cm].
  \item We (drive) \trou[2cm] to Thiès when the car (break) \trou[2cm] down.
\end{enumerate}
\end{exo}
\begin{corrige}[3.1]
1. \en{was walking} --- \en{saw} 2. \en{was cooking} --- \en{went}
3. \en{were you doing} 4. \en{wasn't listening} --- \en{didn't understand}
5. \en{were driving} --- \en{broke}
\end{corrige}

\section{Unité 4 --- Le \en{present perfect} : le chapitre le plus important du livre}

\begin{regle}[Formation]
\en{have / has} + \textbf{participe passé} (3\textsuperscript{e} forme) :\\
\en{I have finished. She has gone. We have seen it. They haven't arrived. Have you eaten?}\\
Contractions : \en{I've, you've, he's, she's, we've, they've} ; \en{haven't, hasn't}.
\end{regle}

\subsection{L'idée en une image}
Le \textbf{passé simple} place l'action dans une boîte fermée du passé : le pont est coupé
avec le présent. Le \textbf{\en{present perfect}} est un \textbf{pont} : l'action est passée,
mais son \emph{effet} ou sa \emph{validité} touche le présent.

\begin{tableau}[Les cinq emplois, avec le test de traduction]
\begin{tabularx}{\linewidth}{@{}p{3.2cm}p{4.6cm}X@{}}
\entete{Emploi} & \entete{Exemple} & \entete{Ce que cela veut dire maintenant} \tabularnewline
\hline
\textbf{1. Résultat présent} & \en{I've lost my keys.} & Je ne les ai \emph{toujours} pas --- je suis bloqué maintenant. (\en{I lost my keys} = simple constat passé, je les ai peut-être retrouvées.) \tabularnewline
\textbf{2. Expérience de vie} & \en{I've been to London twice.} & Dans ma vie jusqu'à maintenant. Aucune date n'est donnée. \tabularnewline
\textbf{3. Durée jusqu'à présent} (avec \en{for/since}) & \en{I've lived here for ten years.} & J'y habite \emph{encore}. (\en{I lived here for ten years} = je n'y habite plus.) \tabularnewline
\textbf{4. Période non terminée} & \en{I've drunk three coffees today.} & La journée n'est pas finie, je peux en boire un quatrième. \tabularnewline
\textbf{5. Nouvelle récente} & \en{The president has arrived.} & On vient de me l'annoncer, c'est d'actualité. \tabularnewline
\end{tabularx}
\end{tableau}

\subsection{Les six mots-signaux}
\begin{tableau}[\en{just, already, yet, ever, never, for, since, recently, so far}]
\begin{tabularx}{\linewidth}{@{}p{2.6cm}p{4.4cm}X@{}}
\entete{Mot} & \entete{Position} & \entete{Exemple et traduction} \tabularnewline
\hline
\en{just} & entre \en{have} et le participe & \en{I've just finished.} --- \textit{Je viens de finir.} \tabularnewline
\en{already} & entre \en{have} et le participe & \en{She's already left.} --- \textit{Elle est déjà partie.} (affirmatif) \tabularnewline
\en{yet} & \textbf{en fin} de phrase & \en{Have you finished yet?} / \en{I haven't eaten yet.} --- \textit{déjà / encore} (question et négation) \tabularnewline
\en{ever} & dans les questions & \en{Have you ever eaten sushi?} --- \textit{As-tu déjà mangé des sushis ?} \tabularnewline
\en{never} & avant le participe & \en{I've never seen snow.} --- \textit{Je n'ai jamais vu la neige.} \tabularnewline
\en{for} & + \textbf{durée} & \en{for two hours, for ten years, for a long time} \tabularnewline
\en{since} & + \textbf{point de départ} & \en{since 2015, since Monday, since I was a child} \tabularnewline
\en{recently / lately} & en fin ou avant le participe & \en{Have you seen him recently?} --- \textit{récemment} \tabularnewline
\en{so far / up to now} & en début ou fin & \en{So far, everything is fine.} --- \textit{Jusqu'ici, tout va bien.} \tabularnewline
\end{tabularx}
\end{tableau}

\begin{piege}[\en{for} et \en{since} : la faute la plus fréquente]
\en{*I live here since 2010.} \ko \ra \en{I've lived here since 2010.} \ok\\
\en{*I am knowing him for five years.} \ko \ra \en{I've known him for five years.} \ok\\
Le français utilise le présent (« j'habite ici depuis 2010 ») ; l'anglais exige le
\en{present perfect}. C'est un automatisme à installer : dès que vous voyez
« depuis », votre cerveau doit déclencher \en{have}.
\end{piege}

\begin{piege}[\en{been} contre \en{gone}]
\en{She has \textbf{been} to Paris.} = elle y est allée (et elle est revenue --- expérience).\\
\en{She has \textbf{gone} to Paris.} = elle y est allée (et elle y est encore --- elle est absente).\\
\en{He's been out.} = il était sorti (il est rentré). \quad \en{He's gone out.} = il est sorti (il n'est pas là).
\end{piege}

\subsection{\en{Present perfect} contre \en{past simple} : le test décisif}
\begin{tableau}[Le seul critère qui compte : la période est-elle finie ?]
\begin{tabularx}{\linewidth}{@{}p{6.0cm}X@{}}
\entete{\en{past simple} (période terminée, date donnée)} & \entete{\en{present perfect} (période ouverte, pas de date)} \tabularnewline
\hline
\en{I saw him \textbf{yesterday}.} & \en{I've seen him \textbf{before}.} \tabularnewline
\en{She \textbf{lived} in Dakar for five years.} (elle n'y vit plus) & \en{She \textbf{has lived} in Dakar for five years.} (elle y vit encore) \tabularnewline
\en{We \textbf{had} breakfast at 7.} & \en{We \textbf{haven't had} breakfast \textbf{yet}.} \tabularnewline
\en{\textbf{Did} you \textbf{go} out last night?} & \en{\textbf{Have} you ever \textbf{been} to Paris?} \tabularnewline
\en{Shakespeare \textbf{wrote} 37 plays.} (il est mort) & \en{My favourite author \textbf{has written} 12 novels.} (il vit encore) \tabularnewline
\en{I \textbf{lost} my keys, but I found them.} & \en{I \textbf{have lost} my keys --- I can't get in.} \tabularnewline
\end{tabularx}
\end{tableau}

\begin{astuce}[L'astuce de la « conversation qui démarre »]
En anglais réel, une conversation commence souvent par un \en{present perfect}, puis
bascule au passé simple dès qu'on entre dans les détails :\\
\en{``Have you seen the new film?''} --- \en{``Yes, I have!''} --- \en{``When did you see it?''}
--- \en{``I saw it on Saturday. It was amazing.''}\\
C'est le schéma naturel : \textbf{ouverture au \en{perfect}, détails au passé simple}.
\end{astuce}

\subsection{Le \en{present perfect continuous}}
\begin{regle}
\en{have / has been} + \en{-ing}. Sert à mettre l'accent sur la \textbf{durée} ou sur
l'\textbf{activité} (souvent encore en cours ou tout juste arrêtée, avec un effet visible).
\end{regle}
\ex{I've been waiting for two hours.}{J'attends depuis deux heures (et j'attends encore).}
\ex{You look tired. --- I've been working all night.}{Tu as l'air fatigué. --- J'ai travaillé toute la nuit.}
\ex{It has been raining. The ground is wet.}{Il a plu (on voit le résultat).}
\ex{How long have you been learning English?}{Depuis combien de temps apprends-tu l'anglais ?}

\begin{tableau}[\en{perfect simple} contre \en{perfect continuous}]
\begin{tabularx}{\linewidth}{@{}p{6.0cm}X@{}}
\entete{Simple : résultat, quantité, achèvement} & \entete{Continu : durée, activité, processus} \tabularnewline
\hline
\en{I've read three books this month.} (quantité achevée) & \en{I've been reading a lot this month.} (activité) \tabularnewline
\en{She's written the report.} (c'est fini) & \en{She's been writing the report all day.} (elle y travaille encore) \tabularnewline
\en{I've known him for years.} (verbe d'état --- jamais de continu) & \en{I've been knowing him} \ko \textbf{interdit} \tabularnewline
\end{tabularx}
\end{tableau}

\begin{piege}[Verbes d'état : jamais au \en{-ing}, même au \en{perfect continuous}]
\en{*I have been knowing him for years.} \ra \en{I have known him for years.}\\
Idem : \en{*have been wanting, *have been believing, *have been owning, *have been meaning}.
\end{piege}

\begin{exo}[4.1 --- \en{present perfect} ou \en{past simple} ?]
\begin{enumerate}
  \item I (lose) \trou[2cm] my phone yesterday.
  \item I (lose) \trou[2cm] my phone --- can you help me find it?
  \item She (visit) \trou[2cm] Morocco three times in her life.
  \item She (visit) \trou[2cm] Morocco in 2019.
  \item We (know) \trou[2cm] each other since 2012.
  \item I (see) \trou[2cm] that film last week.
  \item (you / ever / eat) \trou[3cm] sushi?
  \item He just (arrive) \trou[2cm].
  \item My grandmother (write) \trou[2cm] six books before she died.
  \item I (not / finish) \trou[2.5cm] my homework yet.
\end{enumerate}
\end{exo}
\begin{corrige}[4.1]
1. \en{lost} 2. \en{have lost} 3. \en{has visited} 4. \en{visited} 5. \en{have known}
6. \en{saw} 7. \en{Have you ever eaten} 8. \en{has just arrived} 9. \en{wrote}
10. \en{haven't finished}
\end{corrige}

\begin{exo}[4.2 --- \en{for} ou \en{since} ?]
\begin{enumerate}
  \item \trou[1.5cm] two years \item \trou[1.5cm] 2023 \item \trou[1.5cm] I was a child
  \item \trou[1.5cm] three hours \item \trou[1.5cm] last Monday \item \trou[1.5cm] a long time
  \item \trou[1.5cm] we met \item \trou[1.5cm] ten minutes
\end{enumerate}
\end{exo}
\begin{corrige}[4.2]
\en{for} : 1, 4, 6, 8 (durées). \quad \en{since} : 2, 3, 5, 7 (points de départ).
\end{corrige}

\begin{exo}[4.3 --- Traduisez (pièges)]
\begin{enumerate}
  \item J'habite ici depuis dix ans. \item Je connais Fatou depuis l'école.
  \item Elle est allée au marché (elle n'est pas là). \item Elle est allée au marché trois fois cette semaine.
  \item J'attends depuis une heure. \item Tu as déjà mangé ? \item Je n'ai jamais vu ça.
  \item Nous avons acheté cette voiture en 2020. \item Il vient de sortir.
  \item Depuis combien de temps apprenez-vous l'anglais ?
\end{enumerate}
\end{exo}
\begin{corrige}[4.3]
1. \en{I've lived here for ten years.} 2. \en{I've known Fatou since school.}
3. \en{She's gone to the market.} 4. \en{She's been to the market three times this week.}
5. \en{I've been waiting for an hour.} 6. \en{Have you eaten yet?} / \en{Have you already eaten?}
7. \en{I've never seen that.} 8. \en{We bought this car in 2020.}
9. \en{He's just gone out.} 10. \en{How long have you been learning English?}
\end{corrige}

\section{Unité 5 --- Le \en{past perfect} : le passé du passé}

\begin{regle}
\en{had} + participe passé (toutes les personnes) : \en{I had eaten, she had left,
they hadn't come, had you seen it?}
\end{regle}
\textbf{Emploi :} marquer qu'une action est \textbf{antérieure} à une autre action passée.
\ex{When I arrived, the train had left.}{Quand je suis arrivé, le train était (déjà) parti.}
\ex{She was tired because she hadn't slept.}{Elle était fatiguée parce qu'elle n'avait pas dormi.}
\ex{I realised I had met him before.}{J'ai réalisé que je l'avais déjà rencontré.}
\ex{He said he had finished.}{Il a dit qu'il avait fini. (discours rapporté)}

\begin{astuce}[Quand il est facultatif]
Si l'ordre chronologique est clair grâce à \en{after} ou \en{before}, le passé simple suffit :
\en{After she finished, she left.} Mais avec \en{when}, le \en{past perfect} lève toute
ambiguïté : \en{When she had finished, she left.} (elle a fini, \emph{puis} elle est partie)
contre \en{When she finished, I left.} (au moment où elle a fini, je suis parti).
\end{astuce}

\section{Unité 6 --- Le futur : quatre façons, quatre nuances}

\begin{tableau}[Les quatre futurs et leurs nuances]
\begin{tabularx}{\linewidth}{@{}p{3.0cm}p{4.0cm}X@{}}
\entete{Forme} & \entete{Emploi} & \entete{Exemples} \tabularnewline
\hline
\en{will} + base & Décision prise \textbf{à l'instant}, promesse, prédiction, offre & \en{It's cold. --- I'll close the window.} / \en{I'll help you.} / \en{It will rain tomorrow.} \tabularnewline
\en{be going to} + base & Intention \textbf{déjà formée}, prévision fondée sur un signe visible & \en{I'm going to study tonight.} / \en{Look at those clouds --- it's going to rain.} \tabularnewline
\en{présent continu} & Plan \textbf{organisé, fixé} (rendez-vous, billet acheté) & \en{I'm meeting Ali at 5.} / \en{We're flying to Paris on Monday.} \tabularnewline
\en{présent simple} & Horaire officiel, programme & \en{The train leaves at 7 am.} / \en{The meeting starts at 9.} \tabularnewline
\end{tabularx}
\end{tableau}

\begin{piege}[Ne traduisez pas « aller + infinitif » systématiquement par \en{going to}]
« Je vais le faire » (décision immédiate) = \en{I'll do it.}\\
« Je vais le faire » (projet décidé avant) = \en{I'm going to do it.}\\
« Je le fais demain » (rendez-vous fixé) = \en{I'm doing it tomorrow.}\\
En conversation, les natifs disent très souvent \en{gonna} : \en{I'm gonna call you}.
\end{piege}

\subsection{\en{will} : les formes contractées et les réponses courtes}
\en{I'll, you'll, he'll, she'll, we'll, they'll} \pron{aï-l, iou-l, hi-l\ldots}\\
Négation : \en{won't} \pron{ouount} --- \en{I won't be late.}\\
Question : \en{Will you come?} \quad Réponse courte : \en{Yes, I will.} / \en{No, I won't.}\\
\en{shall} (surtout britannique, avec \en{I/we}) : \en{Shall we go?} (On y va ?),
\en{Shall I open the window?} (Je t'ouvre la fenêtre ?)

\subsection{Le futur continu et le futur antérieur (niveau B1-B2)}
\ex{This time tomorrow, I'll be flying to Dakar.}{Demain à cette heure-ci, je serai en train de voler vers Dakar.}
\ex{By 2028, I will have finished my degree.}{D'ici 2028, j'aurai terminé mes études.}
\ex{Will you have left by 6 pm?}{Seras-tu parti avant 18 h ?}

\subsection{Les subordonnées de temps : présent, pas futur !}
\begin{regle}
Après \en{when, as soon as, before, after, until, if, while}, on emploie le
\textbf{présent} pour parler de l'avenir.
\end{regle}
\ex{I'll call you when I \textbf{arrive}.}{Je t'appellerai quand j'arriverai. (\en{*when I will arrive})}
\ex{We'll eat as soon as he \textbf{comes}.}{Nous mangerons dès qu'il viendra.}
\ex{If it \textbf{rains}, we'll stay home.}{S'il pleut, nous resterons à la maison.}

\begin{exo}[6.1 --- Choisissez le futur]
\begin{enumerate}
  \item The phone is ringing. --- I (answer) \trou[2cm] it.
  \item I (visit) \trou[2.5cm] my uncle next weekend, I've already bought the ticket.
  \item Look out! You (fall) \trou[2cm]!
  \item The conference (begin) \trou[2cm] at 9 am tomorrow.
  \item I promise I (not / tell) \trou[2.5cm] anyone.
  \item We (study) \trou[2.5cm] harder next year, that's our plan.
  \item I'll send it as soon as I (get) \trou[2cm] the file.
  \item This time next week, I (lie) \trou[3cm] on the beach.
\end{enumerate}
\end{exo}
\begin{corrige}[6.1]
1. \en{will answer} / \en{'ll answer} (décision immédiate) 2. \en{am visiting} (plan fixé)
3. \en{are going to fall} (signe visible) 4. \en{begins} (horaire)
5. \en{won't tell} 6. \en{are going to study} (intention) 7. \en{get}
8. \en{will be lying} (futur continu)
\end{corrige}

\section{Unité 7 --- \en{used to} et les habitudes passées}

\begin{regle}
\en{used to} + base = une habitude ou un état \textbf{révolu} du passé.\\
\en{I used to smoke.} (je fumais, maintenant non)\\
Négation : \en{I didn't use to smoke.} \quad Question : \en{Did you use to smoke?}
\end{regle}

\begin{tableau}[Trois formes à ne jamais confondre]
\begin{tabularx}{\linewidth}{@{}p{3.4cm}p{4.0cm}X@{}}
\entete{Forme} & \entete{Sens} & \entete{Exemple} \tabularnewline
\hline
\en{used to} + base & habitude passée révolue & \en{I used to live in Dakar.} --- \textit{J'habitais Dakar (avant).} \tabularnewline
\en{be used to} + nom / \en{-ing} & être habitué à (état actuel) & \en{I'm used to the heat.} --- \textit{Je suis habitué à la chaleur.} \tabularnewline
\en{get used to} + nom / \en{-ing} & s'habituer à (processus) & \en{I'm getting used to waking up early.} --- \textit{Je m'habitue à me lever tôt.} \tabularnewline
\end{tabularx}
\end{tableau}

\begin{piege}[Prononciation et orthographe]
\en{used to} se prononce \pron{iou-s-tou}, jamais \pron{iouz-dou}. Et dans
\en{I'm used to \textbf{working}}, le verbe qui suit est en \en{-ing} --- car \en{to} y est
une préposition, pas un infinitif. \en{*I'm used to work} est une faute classique.
\end{piege}

\ex{When I was young, I used to play football every day.}{Quand j'étais jeune, je jouais au football tous les jours.}
\ex{He didn't use to like coffee, but now he drinks it every morning.}{Il n'aimait pas le café avant, mais maintenant il en boit chaque matin.}
\ex{She would sit by the window for hours.}{Elle s'asseyait près de la fenêtre pendant des heures. (\en{would} pour les habitudes passées répétées --- jamais pour les états : \en{*I would have a car})}

\section{Unité 8 --- Le tableau de bord de tous les temps}
Gardez ce tableau sous les yeux pendant deux mois : il vaut tous les résumés.

\begin{tableau}[Le verbe \en{to work} et le verbe \en{to go} à tous les temps]
\footnotesize
\begin{longtable}{@{}p{3.4cm}p{5.0cm}p{5.4cm}@{}}
\entete{Temps} & \entete{Affirmation} & \entete{Signal / sens} \tabularnewline
\endfirsthead
\entete{Temps} & \entete{Affirmation} & \entete{Signal / sens} \tabularnewline
\endhead
\hline
Présent simple & \en{I work. He works.} & \en{every day, usually, always} \tabularnewline
Présent continu & \en{I am working.} & \en{now, at the moment, today} \tabularnewline
Passé simple & \en{I worked. I went.} & \en{yesterday, last week, in 2010, ago} \tabularnewline
Passé continu & \en{I was working. I was going.} & \en{while, when, at 8 pm yesterday} \tabularnewline
\en{Present perfect} & \en{I have worked. I have gone.} & \en{ever, never, just, already, yet, for, since, so far} \tabularnewline
\en{Perf. continuous} & \en{I have been working.} & \en{for two hours, all day, how long, since} \tabularnewline
\en{Past perfect} & \en{I had worked. I had gone.} & \en{before, after, by the time, when} \tabularnewline
\en{Past perf. continuous} & \en{I had been working.} & durée avant un moment passé \tabularnewline
Futur \en{will} & \en{I will work. I'll go.} & décision, promesse, \en{tomorrow, next year} \tabularnewline
Futur \en{going to} & \en{I am going to work.} & intention, plan, évidence \tabularnewline
Futur continu & \en{I will be working.} & \en{this time tomorrow} \tabularnewline
Futur antérieur & \en{I will have worked.} & \en{by 2030, by the time} \tabularnewline
\end{longtable}
\end{tableau}

\begin{exo}[8.1 --- Le grand test des temps (20 phrases)]
Conjuguez correctement :
\begin{enumerate}
  \item Right now, I (read) \trou[2cm] a book about Africa.
  \item I (read) \trou[2cm] three books this year.
  \item I (read) \trou[2cm] that book last year.
  \item She (live) \trou[2cm] in Paris since 2018.
  \item She (live) \trou[2cm] in Paris from 2010 to 2015.
  \item We (wait) \trou[2.5cm] for you for an hour!
  \item When I (arrive) \trou[2cm], the meeting (already / start) \trou[2.5cm].
  \item Look at the sky --- it (rain) \trou[2.5cm].
  \item I (call) \trou[2cm] you as soon as I (finish) \trou[2cm].
  \item Every morning he (drink) \trou[2cm] tea, but yesterday he (drink) \trou[2cm] coffee.
  \item (you / ever / be) \trou[3cm] to the United States?
  \item I (not / see) \trou[2.5cm] him since Monday.
  \item By next June, they (build) \trou[3cm] the whole house.
  \item While I (walk) \trou[2.5cm] home, I (meet) \trou[2cm] an old friend.
  \item He (be) \trou[2cm] tired because he (work) \trou[3cm] all night.
  \item This time next year, I (study) \trou[3cm] at university.
  \item I (used to / play) \trou[2.5cm] the guitar when I was a teenager.
  \item She (get) \trou[2.5cm] used to the noise now, it's fine.
  \item The shop (open) \trou[2cm] at 9 and (close) \trou[2cm] at 6 every day.
  \item (you / watch) \trou[3cm] TV when I called?
\end{enumerate}
\end{exo}
\begin{corrige}[8.1]
1. \en{am reading} 2. \en{have read} 3. \en{read} \pron{rèd} 4. \en{has lived}
5. \en{lived} 6. \en{have been waiting} 7. \en{arrived} --- \en{had already started}
8. \en{is going to rain} 9. \en{will call / 'll call} --- \en{finish}
10. \en{drinks} --- \en{drank} 11. \en{Have you ever been} 12. \en{haven't seen}
13. \en{will have built} 14. \en{was walking} --- \en{met} 15. \en{was} ---
\en{had been working} 16. \en{will be studying} 17. \en{used to play}
18. \en{is getting} 19. \en{opens} --- \en{closes} 20. \en{Were you watching}
\end{corrige}

\begin{mission}[Chapitre 9 --- les trois récits]
Enregistrez-vous trois fois, une fois par jour :\\
\textbf{Jour 1 --- hier :} 2 minutes au passé simple et au passé continu
(\en{Yesterday I woke up at\ldots\ While I was\ldots, I\ldots}).\\
\textbf{Jour 2 --- votre vie :} 2 minutes au \en{present perfect}
(\en{I have lived\ldots\ I have never\ldots\ I have known\ldots for\ldots}).\\
\textbf{Jour 3 --- l'avenir :} 2 minutes avec \en{will}, \en{going to} et le présent continu
(\en{Next year I'm going to\ldots\ In ten years I'll\ldots\ On Saturday I'm meeting\ldots}).\\
Réécoutez le jour 4 : comptez vos erreurs de temps. Ce sont vos priorités de la semaine
suivante.
\end{mission}
